\subsection{Ideal calculus over arbitrary affine spectra}
\label{subsec:general-ideal-functions}

Let $A$ be any commutative ring and let $x$ be a real multiplicative
ultrametric seminorm on $A$. These statements require neither a base field
nor a prescribed absolute value. They specialize to relative analytifications
over arbitrary valued fields by restriction to their relative spectra. The
finite generator formulas impose boundedness on the coefficient ring; for
nontrivially valued bases this ring is typically a ring of integral elements,
and boundedness is not asserted for the entire base field.

\begin{definition}[Bounded spectra and intrinsic ideal functions]
\label{def:general-ideal-norm}
\lean{EnvelopeProperty.Berkovich.Spectrum.IsBounded,
EnvelopeProperty.Berkovich.Spectrum.idealNorm,
EnvelopeProperty.Berkovich.Spectrum.idealLog,
EnvelopeProperty.Berkovich.Spectrum.idealNorm_le_one,
EnvelopeProperty.Berkovich.Spectrum.idealLog_le_zero}
\leanok
\uses{spectrum:affine-seminorm,spectrum:topology}
\source{fhx-envelope-20261028:subsec:prelim-berkovich}
The point $x$ is bounded if $|a|_x\le1$ for all $a\in A$. For an ideal
$I\subseteq A$, define
\[
 |I|_x=\sup_{a\in I}|a|_x,\qquad \log|I|(x)=\log|I|_x.
\]
The ideal norm has values in $[0,+\infty]$ and its logarithm in the
extended real line. These are intrinsic functions; the centered local ideal
functions of \cite[\S2.1]{FHX} are obtained on bounded coefficient rings.
We use $0\cdot(+\infty)=0$ for nonnegative extended numbers and
$(-\infty)+(+\infty)=-\infty$ for the mixed infinite sum, so the
extended logarithm transforms every product into a sum. On bounded
points ideal norms never equal $+\infty$, so the mixed infinite case
does not occur for centered ideal functions.
\end{definition}

\begin{lemma}[Bounded generators and vanishing]
\label{lem:general-ideal-norm-generators}
\lean{EnvelopeProperty.Berkovich.Spectrum.idealNorm_span,
EnvelopeProperty.Berkovich.Spectrum.idealNorm_span_finset,
EnvelopeProperty.Berkovich.Spectrum.idealNorm_span_finset_eq_of_span_eq,
EnvelopeProperty.Berkovich.Spectrum.idealNorm_eq_zero_iff,
EnvelopeProperty.Berkovich.Spectrum.idealLog_eq_bot_iff}
\leanok
\uses{def:general-ideal-norm}
For a bounded point $x$ and $S\subseteq A$,
$|(S)|_x=\sup_{a\in S}|a|_x$. For finite $S$ the supremum is the maximum,
independent of the generating family; the empty maximum is zero.
At every point, bounded or not, $|I|_x=0$ if and only if
$I\subseteq\operatorname{supp}(x)$, equivalently $\log|I|(x)=-\infty$.
\end{lemma}
\begin{proof}
Set $M=\sup_{a\in S}|a|_x$. The elements with seminorm at most $M$
form an ideal: zero has seminorm zero, sums preserve the bound by the
ultrametric inequality, and multiplication by any coefficient preserves the
bound because $x$ is bounded. This ideal contains $S$ and hence $(S)$.
The reverse inequality follows from $S\subseteq(S)$. For a finite family
the supremum is its maximum. The supremum over an ideal is zero precisely
when every ideal element has seminorm zero, which is the asserted support
criterion. The extended logarithm equals $-\infty$ precisely at zero.
\end{proof}

\begin{lemma}[Intrinsic ideal arithmetic]
\label{lem:general-ideal-norm-arithmetic}
\lean{EnvelopeProperty.Berkovich.Spectrum.idealNorm_sup,
EnvelopeProperty.Berkovich.Spectrum.idealNorm_mul,
EnvelopeProperty.Berkovich.Spectrum.idealNorm_pow_succ,
EnvelopeProperty.Berkovich.Spectrum.idealNorm_pow}
\leanok
\uses{def:general-ideal-norm}
At every point $x$, without a boundedness assumption,
\[
 |I+J|_x=\max\{|I|_x,|J|_x\},\qquad |IJ|_x=|I|_x|J|_x.
\]
For $n\ge1$, $|I^n|_x=|I|_x^n$. If $x$ is bounded, the unit ideal
has norm one and this power identity holds also for $n=0$.
\end{lemma}
\begin{proof}
Every element of $I+J$ is $a+b$ with $a\in I$, $b\in J$; the
ultrametric inequality gives the upper bound, while the inclusions of
$I$ and $J$ give the lower bound. To prove the product upper bound set
$M=|I|_x|J|_x$ and
$K=\{c\in A: |rc|_x\le M\text{ for every }r\in A\}$.
The ultrametric inequality makes $K$ closed under addition, it contains
zero, and replacing $r$ by $rs$ proves closure under multiplication by
$s\in A$. Thus $K$ is an ideal. For $a\in I$, $b\in J$, and any $r$,
$ra\in I$ and
$|r(ab)|_x=|ra|_x|b|_x\le |I|_x|J|_x$, so $ab\in K$.
Hence $IJ\subseteq K$. Taking $r=1$ bounds every element of $IJ$ by
$M$. Conversely, every product $ab$ belongs to $IJ$; taking suprema
and using preservation of suprema by multiplication in $[0,+\infty]$
gives the opposite inequality. Induction starting at $n=1$ proves the
positive power formula. For a bounded point the norm of the unit ideal
is at most one, and its element $1$ gives the reverse bound; this adds
the case $n=0$.
\end{proof}

\begin{lemma}[Generic logarithmic ideal functions]
\label{lem:general-ideal-log}
\lean{EnvelopeProperty.Berkovich.Spectrum.idealLog_sup,
EnvelopeProperty.Berkovich.Spectrum.idealLog_mul,
EnvelopeProperty.Berkovich.Spectrum.idealLog_pow_succ,
EnvelopeProperty.Berkovich.Spectrum.idealLog_pow,
EnvelopeProperty.Berkovich.Spectrum.idealLog_span_finset,
EnvelopeProperty.Berkovich.Spectrum.continuousOn_idealLog_of_span_eq}
\leanok
\uses{lem:general-ideal-norm-generators,lem:general-ideal-norm-arithmetic,
spectrum:logarithmic-evaluation}
At every point, $\log|I+J|$ is the maximum of $\log|I|$ and
$\log|J|$, and $\log|IJ|=\log|I|+\log|J|$. The positive power
formulas also hold without boundedness. At bounded points the power
formula includes exponent zero, and a finitely generated ideal has
logarithmic value equal to the finite maximum of the logarithms of its
generators, with the empty maximum equal to $-\infty$. Its logarithmic
function is continuous on the bounded locus.
\end{lemma}
\begin{proof}
The extended logarithm preserves maxima and transforms multiplication
into addition, including the zero and infinite values. Applying it to
the norm identities proves the arithmetic formulas. On bounded points
it transforms the finite generator maximum into the maximum of the
logarithms. Evaluation on a fixed ring element is continuous in the
spectrum topology, as is the extended logarithm, so this finite maximum
is continuous. For no generators the function is the continuous constant
$-\infty$.
\end{proof}

\begin{lemma}[Generic affine pullbacks]
\label{lem:general-ideal-pullback}
\lean{EnvelopeProperty.Berkovich.Spectrum.IsBounded.comap,
EnvelopeProperty.Berkovich.Spectrum.idealNorm_le_map,
EnvelopeProperty.Berkovich.Spectrum.idealNorm_map,
EnvelopeProperty.Berkovich.Spectrum.idealLog_map}
\leanok
\uses{lem:general-ideal-norm-generators,spectrum:pullback}
Let $\varphi\colon A\to B$ be any homomorphism of commutative rings,
$y\in\operatorname{Spectrum}(B)$ and $x=\varphi^*(y)$. Then
$|I|_x\le|IB|_y$. If $y$ is bounded, so is $x$, and equality holds;
in that case $\log|IB|(y)=\log|I|(x)$.
\end{lemma}
\begin{proof}
For $a\in I$, $\varphi(a)\in IB$ and
$|a|_x=|\varphi(a)|_y$, giving the inequality. If $y$ is bounded then
$|\varphi(a)|_y\le1$ for all $a\in A$, proving boundedness of $x$.
The extended ideal $IB$ is generated by the images of $I$. The generator
identity at $y$ therefore identifies its norm with
$\sup_{a\in I}|\varphi(a)|_y=|I|_x$. Taking logarithms gives the
last assertion.
\end{proof}
